% TOP_PROBLEM: {"id": 1275, "title": "Hilbert-Pólya Conjecture", "category_id": 1, "category": {"id": 1, "name": "number_theory", "display_name": "Number Theory", "description": "Properties of integers, prime numbers, Diophantine equations.", "slug": "number-theory", "order_index": 1, "created_at": "2026-07-31T15:26:25.670Z"}, "status": "open"}
% FIRST_POSTED: null
% ENTRY_KIND: statement_only
% Imported from: batch05.tex; UnsolvedMath ID 1275
\documentclass[11pt]{article}
\usepackage[margin=25mm]{geometry}
\usepackage{fontspec}
\setmainfont{Latin Modern Roman}
\usepackage{amsmath,amssymb,mathtools}
\usepackage{unicode-math}
\setmathfont{Latin Modern Math}
\usepackage{xurl,hyperref}
\hypersetup{colorlinks=true,urlcolor=blue}
\setcounter{secnumdepth}{0}
\setlength{\parindent}{0pt}
\setlength{\parskip}{5pt}
\setlength{\emergencystretch}{3em}
\newcommand{\sref}[2]{\hyperlink{source:#1:#2}{[S#2]}}
\newcommand{\cataloguescope}[1]{\paragraph{Recorded target.}#1}
\begin{document}
% UnsolvedMath ID 1275; source catalogue code NT-068
\setcounter{section}{1274}
\section{Hilbert-Pólya Conjecture}\label{1275}
{\small\sffamily UnsolvedMath ID 1275 · Number Theory}
\subsection{Definitions and mathematical statement}

\paragraph{Shared notation.}$\mathbb N=\{1,2,\ldots\}$, and $\mathbb Z,\mathbb Q,\mathbb R,\mathbb C$ denote the integers, rationals, reals and complex numbers. Any other conventions are specified locally.


Let $\zeta(s)$ be the meromorphic continuation of $\sum_{n=1}^{\infty}n^{-s}$, initially defined for $\operatorname{Re}s>1$. A nontrivial zero is a zero $\rho$ with $0<\operatorname{Re}\rho<1$, counted with its analytic multiplicity. A self-adjoint operator $A$ on a complex Hilbert space $\mathcal H$ is a densely defined linear operator whose domain and action equal those of its adjoint. Here discrete spectrum means isolated real eigenvalues with finite-dimensional eigenspaces, with no finite accumulation point.
\paragraph{Statement recorded in the supplied source.}
Can one construct a complex Hilbert space $\mathcal H$ and a self-adjoint operator $A$ with purely discrete spectrum such that the multiset of nontrivial zeros of $\zeta$ is exactly
\[
\{\tfrac12+i\lambda:\lambda\text{ is an eigenvalue of }A\},
\]
where each eigenvalue is repeated according to the dimension of its eigenspace? \sref{1275}{2}

\subsection{Short English statement}
Can the nontrivial zeta zeros be realized, with their multiplicities, by the real eigenvalues of a self-adjoint operator after shifting by one-half and multiplying by the imaginary unit?
\subsection{Sources}
\begin{description}
\item[\hypertarget{source:1275:1}{S1}] ULAM.AI and the UnsolvedMath Contributors, UnsolvedMath: A Curated Collection of Open Mathematics Problems, record 1275 (NT-068). CC BY 4.0. \url{https://huggingface.co/datasets/ulamai/UnsolvedMath}.
\item[\hypertarget{source:1275:2}{S2}] Sebastian Endres and Frank Steiner, \emph{The Berry--Keating operator on $L^2(\mathbb R_{>},dx)$ and on compact quantum graphs with general self-adjoint realizations}, arXiv:0912.3183v5 (2011), Section 1, pp. 2--3. \url{https://arxiv.org/pdf/0912.3183}.
\end{description}
\label{end:1275}
\end{document}
